% ===========================================================================
\section{Study Area and Data}
\label{sec:data}
% ===========================================================================

\subsection{Study area}

Agriculture in Pakistan is concentrated in the irrigated Indus basin, which extends through
Punjab and Sindh and supports most of the country's crop and livestock production.
Approximately four-fifths of Pakistan's cropland is irrigated through canals and tube wells.
The country also has a large livestock population, with buffalo accounting for a substantial
share, and is among the world's largest consumers of urea fertiliser.

These characteristics shape both the composition and the geography of agricultural
emissions. Enteric fermentation and soil nitrous oxide dominate the greenhouse-gas
inventory, while emissions are spatially concentrated in the irrigated canal belt.

\subsection{Activity data}

National activity data are drawn primarily from official Pakistani sources. Livestock
populations are taken from the 2024 PBS 7th Agricultural Census \citep{pbscensus2024}. The
Pakistan Economic Survey \citep{econsurvey2024ag} is used as a cross-check. The two sources are
not expected to match exactly because the census is based on enumeration, whereas the
Economic Survey includes projections derived from earlier census rounds.
Within the census itself the national summary table and the district tables are separate
transcriptions, and they differ by a handful of head in two species. Neither is edited to force
agreement: the national table is the emissions control, the district tables drive the spatial
allocation, and the residual per species is published with the normalisation factor it implies in
\texttt{outputs/agriculture\_livestock\_reconciliation\_2024.csv}.

Crop area and production data are taken from the Pakistan Economic Survey's Agriculture
chapter for FY2023--24 \citep{econsurvey2024ag} and PBS agricultural statistics.
Synthetic-nitrogen offtake is the audited full-year figure from the National Fertilizer
Development Centre's Fertilizer Review for the same year \citep{nfdc2024}. Both documents are
archived in this repository; the citations name the editions the engine reads, not the most
recent editions published. FAOSTAT activity series are used to reconstruct emissions
for 1990--2023, while the 2024 inventory is anchored to the latest Pakistani census and
administrative data \citep{faostat}.

\subsection{Geospatial data}

Spatial allocation uses the open global datasets listed in Table~\ref{tab:data}.

Livestock emissions are allocated using species-specific densities from Gridded Livestock of
the World~4 \citep{glw4}. Rice-related emissions ride MapSPAM~2020 physical rice
area \citep{mapspam2020}: the raster (V2r2) is held and hash-pinned since the
spatial facility program, its former cropland-fallback class deleted from the
allow-list so a checkout without it refuses the build, and the own-year GCD-Rice
2023 classification is published beside it as a crosscheck and a priced
district-level bracket. Crop-residue
burning is allocated using NASA FIRMS VIIRS
active-fire detections from the reporting period's own burning windows
\citep{firms}. The held record spans all nine fiscal years---\num{1664876} unique
detections fetched over per-crop burning windows derived from the FAO crop calendar's
own harvest dates (rice, wheat, cotton, maize and sugarcane, each harvest followed by a
45-day residue-clearing tail, clipped to its fiscal year)---but the spatial surface
consumes only the reporting period's own dates, \num{166393} in-grid detections inside
FY2023--24. Period-matching is a correction this series already made once (an earlier
extraction spanned October 2024--May 2025, one season after the period, and that offset
was declared until the period-matched extraction replaced it), so the nine-year record
is deliberately not pooled into a climatology here: each year's fires allocate only
that year. The window shapes only the within-province fire geography, never a
magnitude.

The fire layer's construction is stated in full so its declared choices are reviewable.
The product is the science-quality \texttt{VIIRS\_SNPP\_SP} archive, requested in
five-day spans (the area API's maximum) over the national bounding box; detections
are deduplicated on (latitude, longitude, acquisition date).
Each detection is weighted by fire radiative power plus one---FRP approximates biomass
burned, and the offset keeps zero-FRP detections from vanishing. Two filters are
deliberately \emph{not} applied to this spatial surface, and their absence is a
declared choice rather than an
omission: no confidence filter (every detection at any confidence enters the surface)
and no crop mask (the only mask is the national boundary, so non-agricultural fires
inside Pakistan contribute weight where they burn). The monthly timing layer is
stricter: there an attribution ladder masks the bounding box's Indian-Punjab
detections, excludes persistent detections within \num{0.01}\si{\degree} of the kiln
registry's coordinates, and keeps cropland detections only---of the \num{169657}
detections in FY2023--24, \num{103747} fall outside Pakistan, \num{4102} sit near
kilns, \num{9380} off cropland, and \num{52428} are attributed to residue burning; the
unattributed counts are published in the burning crosscheck, never forced into the
profile. A district shape carrying
residue-burning mass in which the layer records no fire falls back to a uniform spread
within that shape, declared per event in the grid fallback ledger; a fetch returning
zero detections refuses to write at all, so an empty fire layer can never look
fetched. Other
cropland-related streams are allocated using the cropland fraction from ESA WorldCover~2021
\citep{worldcover2021}.

Provincial boundaries are obtained from geoBoundaries \citep{runfola2020}, while district
boundaries are based on GADM or GAUL ADM2 polygons. District allocations are evaluated
against crop-area statistics published by the Ministry of National Food Security and
Research \citep{mnfsrcrops}. National emission totals are compared with FAOSTAT, EDGAR and
Climate TRACE.

\begin{table}[!htbp]
\centering
\footnotesize
\begin{tabularx}{\linewidth}{@{}llX@{}}
\toprule
\textbf{Dataset} & \textbf{Use} & \textbf{Role} \\
\midrule
PBS 7th Ag.\ Census & activity & livestock, crop area \\
Economic Survey & activity & crops, livestock cross-check \\
NFDC & activity & synthetic nitrogen \\
FAOSTAT (QCL/RFN/GT) & activity/bench. & time series, benchmark \\
EDGAR 2024/2025 & benchmark & independent national totals \\
Climate TRACE v7.1.9 & benchmark & independent national totals \\
MNFSR district crops & records & district validation (rice, wheat) \\
\GLW{} (2020) & proxy & livestock density \\
MapSPAM 2020 & proxy (held, hash-pinned) & rice area \\
ESA WorldCover 2021 & proxy & cropland fraction \\
NASA FIRMS (VIIRS) & proxy & fire occurrence \\
geoBoundaries / GADM & boundary & province / ADM2 district polygons \\
\bottomrule
\end{tabularx}
\caption{Principal datasets used in the inventory. Pakistani official statistics provide
the national and administrative activity data. Open global products provide spatial proxies
and external benchmarks. Inputs used only in the exploratory land-sector extension are
listed in Supplementary Note~S1.}
\label{tab:data}
\end{table}

\subsection{Status of the estimates}

The estimates remain provisional wherever this inventory holds no nationally measured
parameter. The
livestock population is anchored to the 2024 agricultural census, but the dairy-cattle share
is derived from FAOSTAT milk-animal counts. Rice harvested area and several agronomic
fractions are based on published literature or default assumptions pending national survey
data. Each such parameter is identified in the input files and carried through the resulting
outputs and uncertainty analysis.
